\begin{derivation}[从\eqnrefs{eq:pinem_fourier_expansion}到\eqnrefs{eq:pinem-energy-variance-dp83}]
令
\begin{equation*}
 x\equiv2|\beta|.
\end{equation*}
正文纯相位散射因子具有 Fourier 展开
\begin{equation*}
 \ee^{\ii x\sin\phi}
 =\sum_{\ell=-\infty}^{\infty}
 J_\ell(x)\ee^{\ii\ell\phi}.
\end{equation*}
对于理想单能入射，输出第 $\ell$ 条边带的振幅只差一个整体相位，因此
\begin{equation*}
 P_\ell=J_\ell^2(x).
\end{equation*}
为了同时获得归一化和所有矩，不逐项求和 $\sum_\ell \ell^nJ_\ell^2(x)$，而构造概率特征函数
\begin{equation*}
 \chi_\ell(\vartheta)
 =\sum_{\ell=-\infty}^{\infty}
 P_\ell \ee^{\ii\ell\vartheta}.
\end{equation*}

考虑同一个相位因子的角向自相关
\begin{equation*}
 C(\vartheta)
 \equiv\frac1{2\pi}
 \int_0^{2\pi}\dd\phi\,
 \ee^{\ii x\sin(\phi+\vartheta)}
 \ee^{-\ii x\sin\phi}.
\end{equation*}
先用 Fourier 展开计算。第一因子为
\begin{equation*}
 \ee^{\ii x\sin(\phi+\vartheta)}
 =\sum_{\ell}J_\ell(x)
 \ee^{\ii\ell\phi}\ee^{\ii\ell\vartheta},
\end{equation*}
第二因子可写成
\begin{equation*}
 \ee^{-\ii x\sin\phi}
 =\sum_mJ_m(x)\ee^{-\ii m\phi},
\end{equation*}
这里利用 $J_m(x)$ 对实 $x$ 为实数。于是
\begin{align*}
 C(\vartheta)
 &=\frac1{2\pi}
 \sum_{\ell,m}J_\ell(x)J_m(x)\ee^{\ii\ell\vartheta}
 \int_0^{2\pi}\dd\phi\,
 \ee^{\ii(\ell-m)\phi}\\
 &=\sum_{\ell,m}J_\ell(x)J_m(x)\ee^{\ii\ell\vartheta}
 \delta_{\ell m}\\
 &=\sum_\ell J_\ell^2(x)\ee^{\ii\ell\vartheta}
 =\chi_\ell(\vartheta).
\end{align*}
因此自相关本身就是边带阶数的特征函数。

另一方面，不使用 Fourier 级数也可以直接算同一个积分。三角恒等式给
\begin{align*}
 \sin(\phi+\vartheta)-\sin\phi
 &=2\sin\frac{\vartheta}{2}
 \cos\!\left(\phi+\frac{\vartheta}{2}\right).
\end{align*}
令
\begin{equation*}
 u=\phi+\frac{\vartheta}{2},
\end{equation*}
由于积分区间跨越完整 $2\pi$ 周期，平移不改变积分：
\begin{align*}
 C(\vartheta)
 &=\frac1{2\pi}\int_0^{2\pi}\dd u\,
 \exp\!\left[
 \ii\,2x\sin\frac{\vartheta}{2}\cos u
 \right].
\end{align*}
Bessel 函数的积分定义
\begin{equation*}
 J_0(y)
 =\frac1{2\pi}\int_0^{2\pi}\dd u\,\ee^{\ii y\cos u}
\end{equation*}
于是给出
\begin{align*}
 \chi_\ell(\vartheta)
 &=J_0\!\left(2x\sin\frac{\vartheta}{2}\right)\\
 &=J_0\!\left(4|\beta|\sin\frac{\vartheta}{2}\right),
\end{align*}
即正文\eqnrefs{eq:pinem-sideband-characteristic-dp83}。令 $\vartheta=0$：
\begin{equation*}
 \sum_\ell P_\ell
 =\chi_\ell(0)=J_0(0)=1,
\end{equation*}
所以边带分布的归一化也同时得到，而不是另行假定。

特征函数的 $n$ 阶导数满足
\begin{equation*}
 \chi_\ell^{(n)}(0)
 =\sum_\ell(\ii\ell)^nP_\ell
 =\ii^n\langle\ell^n\rangle.
\end{equation*}
由于 $P_{-\ell}=P_\ell$，所有奇数矩为零。为了得到二阶和四阶矩，先把复合函数在 $\vartheta=0$ 附近展开。令
\begin{equation*}
 y(\vartheta)
 =4|\beta|\sin\frac{\vartheta}{2}.
\end{equation*}
正弦展开为
\begin{align*}
 y(\vartheta)
 &=4|\beta|
 \left(\frac\vartheta2-\frac1{6}\frac{\vartheta^3}{8}
 +O(\vartheta^5)\right)\\
 &=2|\beta|\vartheta
 -\frac{|\beta|}{12}\vartheta^3
 +O(\vartheta^5).
\end{align*}
因此
\begin{align*}
 y^2
 &=4|\beta|^2\vartheta^2
 -\frac{|\beta|^2}{3}\vartheta^4
 +O(\vartheta^6),\\
 y^4
 &=16|\beta|^4\vartheta^4
 +O(\vartheta^6).
\end{align*}
再用
\begin{equation*}
 J_0(y)
 =1-\frac{y^2}{4}+\frac{y^4}{64}+O(y^6),
\end{equation*}
逐项代入：
\begin{align*}
 \chi_\ell(\vartheta)
 &=1-\frac14
 \left(4|\beta|^2\vartheta^2
 -\frac{|\beta|^2}{3}\vartheta^4\right)
 +\frac1{64}
 16|\beta|^4\vartheta^4
 +O(\vartheta^6)\\
 &=1-|\beta|^2\vartheta^2
 +\left(
 \frac{|\beta|^2}{12}
 +\frac{|\beta|^4}{4}
 \right)\vartheta^4
 +O(\vartheta^6).
\end{align*}
另一方面，按概率矩展开
\begin{align*}
 \chi_\ell(\vartheta)
 &=1
 +\ii\langle\ell\rangle\vartheta
 -\frac{\langle\ell^2\rangle}{2}\vartheta^2
 -\frac{\ii\langle\ell^3\rangle}{6}\vartheta^3
 +\frac{\langle\ell^4\rangle}{24}\vartheta^4
 +O(\vartheta^5).
\end{align*}
比较 $\vartheta^2$ 系数：
\begin{equation*}
 -\frac{\langle\ell^2\rangle}{2}=-|\beta|^2
 \quad\Rightarrow\quad
 \langle\ell^2\rangle=2|\beta|^2.
\end{equation*}
比较 $\vartheta^4$ 系数：
\begin{align*}
 \frac{\langle\ell^4\rangle}{24}
 &=\frac{|\beta|^2}{12}+\frac{|\beta|^4}{4},\\
 \langle\ell^4\rangle
 &=2|\beta|^2+6|\beta|^4.
\end{align*}
这得到正文\eqnrefs{eq:pinem-sideband-moments-dp83}。

最后把边带阶数映射成理想均匀能量格。在该有效模型内
\begin{equation*}
 E_{\rm out}=E_0+\ell\hbar\omega.
\end{equation*}
所以
\begin{align*}
 \langle E_{\rm out}\rangle
 &=E_0+\hbar\omega\langle\ell\rangle
 =E_0,\\
 \operatorname{Var}(E_{\rm out})
 &=\langle(E_{\rm out}-E_0)^2\rangle\\
 &=(\hbar\omega)^2\langle\ell^2\rangle\\
 &=2|\beta|^2(\hbar\omega)^2,
\end{align*}
即正文\eqnrefs{eq:pinem-energy-variance-dp83}。若真实反冲使通道能量不再满足 $\mathcal E_\ell=E_0+\ell\hbar\omega$，最后这一步必须改为直接对真实 $\mathcal E_\ell$ 做加权求和；前面的 Bessel 概率矩只描述理想均匀格模型内部的统计。
\end{derivation}
